% GENERATED FILE. Edit canonical lessons, then run npm run generate:books.
% canonical-source-hash: fnv1a64:edb6972790f1fcb4
% canonical-lessons: ES-C691-agricultor, ES-C691-campesino, ES-C691-finca, ES-C691-vinedo, ES-C691-cultivar

\chapter{La vida en el campo}
\label{ch:es-la-vida-en-el-campo}
\hlchaptermodality{691}

\begin{chapteropening}
\textbf{By the end of this chapter:} \emph{I can talk about life on the land: farmers, a country estate, a vineyard and growing crops.}

Complete the last lesson of chapter 691: I can talk about life on the land: farmers, a country estate, a vineyard and growing crops.
\end{chapteropening}

\section[el agricultor]{el agricultor — the farmer}

\label{lesson:ES-C691-agricultor}



\begin{quote}
Before the new one: say the Spanish for the swimsuit, then the Spanish for to lie down.
\end{quote}

\subsection*{You'll want to know: el agricultor}
\textbf{el agricultor} — \textquotedblleft{}the farmer\textquotedblright{}.

\begin{quote}
El agricultor trabaja desde muy temprano.
\end{quote}

\begin{quote}
\emph{The farmer works from very early.}
\end{quote}

A woman is \emph{la agricultora}. \emph{La agricultura} is farming.

\begin{cousinweb}
\emph{agricultor} is a learned borrowing from Latin \emph{agricultor} \textquotedblleft{}a farmer\textquotedblright{} (EtymDB), made from \emph{ager} \textquotedblleft{}field\textquotedblright{} and \emph{cultor} \textquotedblleft{}one who tills\textquotedblright{}.
\end{cousinweb}

\begin{grammarlens}[title={What you've built}]
One more word for telling, explaining and describing at B1.
\end{grammarlens}

\subsection*{Your turn}
\begin{itemize}
\raggedright
  \item \emph{Say it:} \emph{el agricultor}
  \item \emph{Say it:} \emph{el agricultor}, once more
  \item \emph{Say it:} a sentence of your own with \emph{el agricultor}
  \item \emph{Recall:} say the Spanish for the swimsuit, then the Spanish for to lie down, then say \emph{el agricultor} again
\end{itemize}

\begin{tcolorbox}[breakable,colback=teal!4,colframe=teal!35!black,title={Before you move on}]
What does el agricultor mean? (\textquotedblleft{}The farmer\textquotedblright{}.) Say it once more.
\end{tcolorbox}
\section[el campesino]{el campesino — the country person, the peasant, the farm worker}

\label{lesson:ES-C691-campesino}



\begin{quote}
Before the new one: say the Spanish for to lie down, then the Spanish for the farmer.
\end{quote}

\subsection*{You'll want to know: el campesino}
\textbf{el campesino} — \textquotedblleft{}the country person, the peasant, the farm worker\textquotedblright{}.

\begin{quote}
Los campesinos venden fruta en el mercado.
\end{quote}

\begin{quote}
\emph{The country people sell fruit at the market.}
\end{quote}

As an adjective it means country: \emph{la vida campesina}.

\begin{cousinweb}
\emph{campesino} is built on \emph{campo} \textquotedblleft{}field, countryside\textquotedblright{}, which is inherited from Latin \emph{campus} \textquotedblleft{}a field, a plain\textquotedblright{}.
\end{cousinweb}

\begin{grammarlens}[title={What you've built}]
One more word for telling, explaining and describing at B1.
\end{grammarlens}

\subsection*{Your turn}
\begin{itemize}
\raggedright
  \item \emph{Say it:} \emph{el campesino}
  \item \emph{Say it:} \emph{el campesino}, once more
  \item \emph{Say it:} a sentence of your own with \emph{el campesino}
  \item \emph{Recall:} say the Spanish for to lie down, then the Spanish for the farmer, then say \emph{el campesino} again
\end{itemize}

\begin{tcolorbox}[breakable,colback=teal!4,colframe=teal!35!black,title={Before you move on}]
What does el campesino mean? (\textquotedblleft{}The country person, the peasant, the farm worker\textquotedblright{}.) Say it once more.
\end{tcolorbox}
\section[la finca]{la finca — the country estate, the farm; the property}

\label{lesson:ES-C691-finca}



\begin{quote}
Before the new one: say the Spanish for the farmer, then the Spanish for the country person.
\end{quote}

\subsection*{You'll want to know: la finca}
\textbf{la finca} — \textquotedblleft{}the country estate, the farm; the property\textquotedblright{}.

\begin{quote}
Mis tíos tienen una finca con caballos.
\end{quote}

\begin{quote}
\emph{My aunt and uncle have a farm with horses.}
\end{quote}

In Spain \emph{una finca} is also a building or plot in town.

\begin{cousinweb}
\emph{finca} is built on old \emph{fincar} \textquotedblleft{}to stay put, to settle\textquotedblright{} (EtymDB), which comes from a Vulgar Latin form of \emph{fīgere} \textquotedblleft{}to fix\textquotedblright{}.
\end{cousinweb}

\begin{grammarlens}[title={What you've built}]
One more word for telling, explaining and describing at B1.
\end{grammarlens}

\subsection*{Your turn}
\begin{itemize}
\raggedright
  \item \emph{Say it:} \emph{la finca}
  \item \emph{Say it:} \emph{la finca}, once more
  \item \emph{Say it:} a sentence of your own with \emph{la finca}
  \item \emph{Recall:} say the Spanish for the farmer, then the Spanish for the country person, then say \emph{la finca} again
\end{itemize}

\begin{tcolorbox}[breakable,colback=teal!4,colframe=teal!35!black,title={Before you move on}]
What does la finca mean? (\textquotedblleft{}The country estate, the farm; the property\textquotedblright{}.) Say it once more.
\end{tcolorbox}
\section[el viñedo]{el viñedo — the vineyard}

\label{lesson:ES-C691-vinedo}



\begin{quote}
Before the new one: say the Spanish for the country person, then the Spanish for the country estate.
\end{quote}

\subsection*{You'll want to know: el viñedo}
\textbf{el viñedo} — \textquotedblleft{}the vineyard\textquotedblright{}.

\begin{quote}
En otoño trabajan en el viñedo.
\end{quote}

\begin{quote}
\emph{In autumn they work in the vineyard.}
\end{quote}

Same family as \emph{la viña} and \emph{el vino}.

\begin{cousinweb}
\emph{viñedo} is inherited from Latin \emph{vīnētum} \textquotedblleft{}a vineyard\textquotedblright{} (EtymDB), from \emph{vīnum} \textquotedblleft{}wine\textquotedblright{}.
\end{cousinweb}

\begin{grammarlens}[title={What you've built}]
One more word for telling, explaining and describing at B1.
\end{grammarlens}

\subsection*{Your turn}
\begin{itemize}
\raggedright
  \item \emph{Say it:} \emph{el viñedo}
  \item \emph{Say it:} \emph{el viñedo}, once more
  \item \emph{Say it:} a sentence of your own with \emph{el viñedo}
  \item \emph{Recall:} say the Spanish for the country person, then the Spanish for the country estate, then say \emph{el viñedo} again
\end{itemize}

\begin{tcolorbox}[breakable,colback=teal!4,colframe=teal!35!black,title={Before you move on}]
What does el viñedo mean? (\textquotedblleft{}The vineyard\textquotedblright{}.) Say it once more.
\end{tcolorbox}
\section[cultivar]{cultivar — to grow (crops), to farm; to cultivate}

\label{lesson:ES-C691-cultivar}



\begin{quote}
Before the new one: say the Spanish for the country estate, then the Spanish for the vineyard.
\end{quote}

\subsection*{You'll want to know: cultivar}
\textbf{cultivar} — \textquotedblleft{}to grow (crops), to farm; to cultivate\textquotedblright{}.

\begin{quote}
Aquí cultivan arroz.
\end{quote}

\begin{quote}
\emph{They grow rice here.}
\end{quote}

\emph{El cultivo} is the crop or the growing of it.

\begin{cousinweb}
\emph{cultivar} is a learned borrowing from Medieval Latin \emph{cultīvāre} (EtymDB), built on \emph{colere} \textquotedblleft{}to till, to tend\textquotedblright{}.
\end{cousinweb}

\begin{grammarlens}[title={What you've built}]
One more word for telling, explaining and describing at B1.
\end{grammarlens}

\subsection*{Your turn}
\begin{itemize}
\raggedright
  \item \emph{Say it:} \emph{cultivar}
  \item \emph{Say it:} \emph{cultivar}, once more
  \item \emph{Say it:} a sentence of your own with \emph{cultivar}
  \item \emph{Recall:} say the Spanish for the country estate, then the Spanish for the vineyard, then say \emph{cultivar} again
\end{itemize}

\begin{tcolorbox}[breakable,colback=teal!4,colframe=teal!35!black,title={Before you move on}]
What does cultivar mean? (\textquotedblleft{}To grow (crops), to farm; to cultivate\textquotedblright{}.) Say it once more.
\end{tcolorbox}
